% One read, two call paths, side by side.
\begin{tikzpicture}[font=\sffamily\scriptsize]
\foreach \x/\name/\id in {0/application/app, 3.4/vendor layer/hal,
                          6.8/I2C1 registers/regs, 10.4/sensor/dev}{
  \node[actor, text width=19mm] (\id) at (\x,0) {\name};
  \draw[lifeline] (\id.south) -- ++(0,-7.4);
}

\draw[msg] (0,-0.9) -- node[above]{memory read, 6 bytes} (3.4,-0.9);
\node[activation, minimum height=24mm] at (3.4,-2.5) {};
\draw[msg] (3.4,-1.35) -- ++(0.75,0) -- ++(0,-0.4)
  node[right, xshift=1mm, font=\sffamily\tiny]{check state} -- ++(-0.75,0);
\draw[msg] (3.4,-2.25) -- ++(0.75,0) -- ++(0,-0.4)
  node[right, xshift=1mm, font=\sffamily\tiny]{take the lock} -- ++(-0.75,0);
\draw[msg] (3.4,-3.15) -- ++(0.75,0) -- ++(0,-0.4)
  node[right, xshift=1mm, font=\sffamily\tiny]{deadline from the tick} -- ++(-0.75,0);
\draw[msg] (3.4,-4.2) -- node[above]{configure and start} (6.8,-4.2);
\draw[msg] (6.8,-4.85) -- node[above]{address, then bytes} (10.4,-4.85);
\draw[reply] (10.4,-5.5) -- node[below]{six bytes} (6.8,-5.5);
\draw[reply] (6.8,-6.1) -- node[below]{flags} (3.4,-6.1);
\draw[msg] (3.4,-6.6) -- ++(0.75,0) -- ++(0,-0.4)
  node[right, xshift=1mm, font=\sffamily\tiny]{map errors and release} -- ++(-0.75,0);
\draw[reply] (3.4,-7.2) -- node[below]{status} (0,-7.2);

% ---------------- the register path, drawn as the same diagram with the
% middle lifeline removed
\node[umlnote, text width=58mm, anchor=north west] at (-0.6,-8.4)
  {The register implementation is this diagram with the middle lifeline taken
   out. The application configures and starts the transfer itself and reads the
   bytes back itself. Nothing checks the state, nothing takes a lock, nothing
   holds a deadline and nothing maps a flag onto a return value.};
\node[umlnote, text width=58mm, anchor=north west] at (5.6,-8.4)
  {So the honest question is not whether the layer is worth its cost but which
   of those four services you need. Add them back one at a time to the register
   version and its cost converges on the layer's. That convergence is the
   result this chapter is built to find.};
\end{tikzpicture}
